\begin{algorithm}[H]
\SetAlgoLined
\setstretch{0.92}
\SetAlFnt{\normalsize}
\SetArgSty{textnormal}
\SetAlgoNlRelativeSize{-1}
\SetInd{0.55em}{1.05em}
\SetKwInput{KwIn}{Input}
\SetKwInput{KwInit}{Initialize}
\SetKwInput{KwOut}{Output}
\SetKwFor{For}{for}{do}{end for}
\SetKwIF{If}{ElseIf}{Else}{if}{then}{else if}{else}{end if}
\SetKw{KwRet}{return}
\caption{问题二分角度多初值有界反演算法}
\label{alg:q2-inversion}
\KwIn{$10^\circ$、$15^\circ$实测反射光谱，问题一双光束反射率$R_{\mathrm{th}}$}
\KwInit{拟合窗口$[2500,3300]\,\mathrm{cm}^{-1}$，参数可行域$\Omega_k$，$90$组低维搜索初值}
\For{$k\in\{10^\circ,15^\circ\}$}{
将反射率百分数转为无量纲小数，并截取拟合窗口\;
\textbf{阶段I：低维盆地搜索}\;
固定Sellmeier参数$\boldsymbol\beta_0$，建立候选集$\mathcal C_k\leftarrow\varnothing$\;
\For{每一组初值$\boldsymbol s\in\mathcal S$}{
对$(d,q_N,n_3)$求解有界非线性最小二乘\;
计算$J_k$与RMSE并写入$\mathcal C_k$\;
}
合并近似重复候选，按RMSE升序保留前$8$个不同盆地\;
\textbf{阶段II：完整参数精修}\;
建立精修集合$\mathcal F_k\leftarrow\varnothing$\;
\For{$\boldsymbol c\in\mathcal C_k$的前$8$个候选}{
释放$7$个Sellmeier参数，在完整十参数可行域内重新最小二乘\;
记录厚度、RMSE、MAPE、$R^2$及边界活跃参数并写入$\mathcal F_k$\;
}
取$\boldsymbol p_k^*\leftarrow\arg\min_{\boldsymbol p\in\mathcal F_k}\mathrm{RMSE}(\boldsymbol p)$\;
记录$d_k\leftarrow d(\boldsymbol p_k^*)$\;
}
计算$d_{\mathrm{SiC}}\leftarrow(d_{10^\circ}+d_{15^\circ})/2$\;
\KwRet $d_{\mathrm{SiC}}$及两个角度的最优候选、拟合指标和局部极小解信息\;
\KwOut{SiC外延层厚度$d_{\mathrm{SiC}}$及反演诊断量}
\end{algorithm}